% section: 極限

\section{極限}\label{節：極限}
  本節では,数列や関数の値がある値に近づくときの考え方を確認する.
  後の節でオイラーの公式を導くために,三角関数と自然対数に関する基本的な極限を準備する.

  \subsection{数列の極限}
    数列$\{a_n\}$の番号$n$を大きくしていくとき,\,$a_n$がある値$L$に限りなく近づくなら,\,$a_n$は$L$に収束するといい,
    \[
      \lim_{n\to\infty}a_n=L
    \]
    と書く.
    これは,\,$L$の近くにどれほど狭い範囲を指定しても,\,$n$を十分大きくすればその範囲に$a_n$が入るという意味である.

    例えば,
    \[
      a_n=\frac{1}{n}
    \]
    とすると,\,$n$を大きくするほど$a_n$は$0$に近づくので,
    \[
      \lim_{n\to\infty}\frac{1}{n}=0
    \]
    である.

  \subsection{関数の極限}
    $x$を$a$に近づけたとき,\,$f(x)$が$L$に近づくなら,
    \[
      \lim_{x\to a}f(x)=L
    \]
    と書く.
    このとき,\,$x=a$での関数の値そのものではなく,\,$a$の近くでの値の変化に注目する.

    極限の計算では,既に求めた極限を組み合わせたり,不等式で上下から挟んだりする.
    例えば,
    \[
      a_n\le b_n\le c_n
    \]
    が十分大きな$n$で成り立ち,\,$a_n$と$c_n$が同じ値$L$に収束するなら,\,$b_n$も$L$に収束する.
    これをはさみうちの原理という.

  \subsection{三角関数の基本極限}
    角度は弧度法で表す.
    $0<x<\dfrac{\pi}{2}$のとき,半径$1$の円の扇形と,その扇形に内接・外接する三角形の面積を比べると,
    \[
      \frac{1}{2}\sin x
      <
      \frac{1}{2}x
      <
      \frac{1}{2}\tan x
    \]
    となる.
    したがって,
    \[
      \cos x<\frac{\sin x}{x}<1
    \]
    である.
    $x\to0+$のとき$\cos x\to1$なので,はさみうちの原理より
    \[
      \lim_{x\to0+}\frac{\sin x}{x}=1
    \]
    を得る.
    また,\,$\dfrac{\sin x}{x}$は偶関数であるから,左側から近づけた場合にも同じ極限となる.
    よって,
    \[
      \lim_{x\to0}\frac{\sin x}{x}=1
    \]
    である.

    さらに,\,$\tan x=\dfrac{\sin x}{\cos x}$および$\lim_{x\to0}\cos x=1$より,
    \[
      \lim_{x\to0}\frac{\tan x}{x}
      =
      \lim_{x\to0}
      \left(
        \frac{\sin x}{x}\cdot\frac{1}{\cos x}
      \right)
      =1
    \]
    となる.

  \subsection{自然対数と指数関数の基本極限}
    以下では,\,$\log x$を自然対数$\log_e x$の意味で用いる.
    $x>-1$に対して,
    \[
      \log(1+x)=\int_0^x\frac{1}{1+t}\,dt
    \]
    である.
    $0<|x|<1$のとき,積分区間で
    \[
      \frac{1}{1+|x|}
      \le
      \frac{1}{1+t}
      \le
      \frac{1}{1-|x|}
    \]
    が成り立つ.
    両辺を区間の長さで平均すると,
    \[
      \frac{1}{1+|x|}
      \le
      \frac{\log(1+x)}{x}
      \le
      \frac{1}{1-|x|}
    \]
    となる.
    $x\to0$のとき両端は$1$に近づくので,
    \[
      \lim_{x\to0}\frac{\log(1+x)}{x}=1
    \]
    を得る.

    この極限を用いると,任意の実数$a$について
    \[
      \lim_{n\to\infty}\left(1+\frac{a}{n}\right)^n=e^a
    \]
    が成り立つ.
    $a=0$のときは明らかである.
    $a\ne0$のとき,\,$n$を十分大きくすれば底は正であり,
    \begin{align*}
      \log\left\{\left(1+\frac{a}{n}\right)^n\right\}
        &=n\log\left(1+\frac{a}{n}\right)\\
        &=a\,
          \frac{
            \log\left(1+\frac{a}{n}\right)
          }{
            \frac{a}{n}
          }\\
        &\longrightarrow a.
    \end{align*}
    よって,両辺の指数を考えれば極限は$e^a$となる.
    この極限は,実数の指数関数を
    \[
      e^a=\lim_{n\to\infty}\left(1+\frac{a}{n}\right)^n
    \]
    と表すものである.

    以上の極限を用いて,次節では複素数の指数表示とオイラーの公式を導く.

